\documentclass[12pt]{article}
\usepackage[utf8]{inputenc}
\usepackage[T1]{fontenc}
\usepackage[margin=1in]{geometry}
\usepackage{setspace}
\usepackage{xcolor}
\usepackage[colorlinks=true, urlcolor=blue!60!black, linkcolor=blue!60!black]{hyperref}
\setlength{\parindent}{0em}
\setlength{\parskip}{0.45em}
\renewcommand{\baselinestretch}{1.0}
\begin{document}
\begin{center}
{\Large \textit{DHS v. Thuraissigiam}, 591 U.S. 103 (2020)\par}
\vspace{0.3em}
{\normalsize Case Brief --- One-Page Edition\par}
\end{center}

\textbf{Memory Jogger.} Habeas secures release, not admission; at the border, Congress's procedures are all the process due (\href{https://supreme.justia.com/cases/federal/us/591/19-161/}{591 U.S. 103}).

\textbf{Facts.} Sri Lankan national Vijayakumar Thuraissigiam was caught 25 yards inside the border and placed in expedited removal; an asylum officer, a supervisor, and an immigration judge found no credible fear of persecution. Section 1252(e)(2) bars habeas review of that finding.

\textbf{Procedural History.} The district court dismissed his habeas petition; the Ninth Circuit held the bar unconstitutional.

\textbf{Votes.} Reversed, 7--2. Alito wrote for five; Thomas concurred; Breyer and Ginsburg concurred in the judgment; Sotomayor and Kagan dissented.

\textbf{Holding.} The bar violates neither the Suspension Clause nor due process as applied: the 1789 writ secured release from unlawful detention, and a noncitizen at the threshold of entry has only the procedures Congress provides.

\textbf{Reasoning.} He sought a new asylum hearing, not release; founding-era cases ordered release, and later immigration review rested on statute. A person caught just after crossing has not entered. Breyer would have decided more narrowly; Sotomayor said the Court ``handcuff[ed]'' judicial oversight of executive detention.

\textbf{Examples.} \emph{Same-Side:} a migrant caught near the Rio Grande seeking a new interview. \emph{Opposite-Side:} a long-time resident detained indefinitely seeking release. \emph{Fence-Sitter:} expedited removal of an eighteen-month interior resident.

\textbf{Critique.} Demanding a 1789 analogue is anachronistic, and the undefined threshold-of-entry line risks unreviewed error.

\textbf{Key Quotation.} ``For these reasons, an alien in respondent's position has only those rights regarding admission that Congress has provided by statute.'' --- Justice Alito, Opinion of the Court, 591 U.S. at 140.

\end{document}
